\documentclass[a4paper,fleqn,longmktitle]{cas-sc}

% TeX Live 2024 renamed the expl3 primitive used by CAS 2.3.
\ExplSyntaxOn
\cs_if_exist:NF \vbox_unpack_clear:N
  { \cs_new_eq:NN \vbox_unpack_clear:N \vbox_unpack_drop:N }
\ExplSyntaxOff

\usepackage[numbers,sort&compress]{natbib}
\usepackage[T1]{fontenc}
\usepackage[utf8]{inputenc}
\usepackage{microtype}
\usepackage{amsmath,amssymb,amsthm,mathtools}
\usepackage{booktabs}
\usepackage{array}
\usepackage{enumitem}
\usepackage{xcolor}
\usepackage{hyperref}
\usepackage[nameinlink,noabbrev]{cleveref}

\hypersetup{
  colorlinks=true,
  linkcolor=blue!55!black,
  citecolor=blue!55!black,
  urlcolor=blue!55!black,
  pdftitle={The exact saturated 6- and 7-Sperner numbers},
  pdfauthor={Jiazhi Chen}
}

\newtheorem{theorem}{Theorem}[section]
\newtheorem{proposition}[theorem]{Proposition}
\newtheorem{lemma}[theorem]{Lemma}
\newtheorem{corollary}[theorem]{Corollary}
\theoremstyle{definition}
\newtheorem{definition}[theorem]{Definition}
\theoremstyle{remark}
\newtheorem{remark}[theorem]{Remark}

\newcommand{\powerset}{\mathcal{P}}
\newcommand{\family}{\mathcal{F}}
\newcommand{\smallcores}{\mathcal{S}}
\newcommand{\largecores}{\mathcal{L}}
\newcommand{\sat}{\operatorname{sat}}
\newcommand{\blocker}{\mathcal{B}}
\newcommand{\bits}[1]{\mathtt{#1}}

\begin{document}

\let\WriteBookmarks\relax
\def\floatpagepagefraction{1}
\def\textpagefraction{.001}

\shorttitle{Exact saturated 6- and 7-Sperner numbers}
\shortauthors{J. Chen}

\title[mode=title]{The exact saturated 6- and 7-Sperner numbers}

\author[1]{Jiazhi Chen}[orcid=0009-0008-9790-9881]
\cormark[1]
\ead{lailai0x394@gmail.com}
\credit{Conceptualization, Methodology, Software, Formal analysis,
Investigation, Validation, Writing -- original draft, Writing -- review and editing}

\affiliation[1]{
  organization={Hangzhou No.2 High School of Zhejiang Province},
  addressline={No. 76 Dongxin Avenue, Binjiang District},
  city={Hangzhou},
  postcode={310053},
  state={Zhejiang},
  country={China}}

\cortext[1]{Corresponding author}

\begin{abstract}
We determine the eventual saturated 6- and 7-Sperner numbers.  At six layers,
the exact value is $\sat(6)=30$.  Its lower bound reduces the two middle
canonical layers to the finite-clutter parameter $m(2,3)$, for which we prove
$m(2,3)=9$ on arbitrary finite ground sets.  At seven layers, a stronger
blocker inequality first gives $m(2,4)=12$ and $\sat(7)\geq47$.  Successive
arbitrary-finite reductions then exclude the equality profiles of sizes 52
and 53 and reduce size 54 to two profiles.  Exact branch arguments and
independently replayed finite kernel enumerations exclude the Fano profile and
all middle splits $(5,11)$, $(6,10)$, $(7,9)$, and $(8,8)$, together with
their blocker duals.  Hence $\sat(7)\geq55$.  An explicit 55-member family
over an eight-element core and a disjoint homogeneous atom gives equality.
Its small/large split is $28/27$, so composition also yields
\[
  \sat(5j+2+s)\leq 2^s(28^j+27^j)
  \qquad(j\geq1,\ 0\leq s\leq4).
\]
Both eventual equalities are also formalized in Lean 4 as attainment-and-
optimality statements on every sufficiently large labelled ground set.  A
certificate-checked encoding separately proves that 56 is optimal only within
the analogous seven-core layered template class.
\end{abstract}

\begin{highlights}
\item The eventual saturated 6-Sperner number is 30.
\item The eventual saturated 7-Sperner number is 55.
\item Arbitrary-finite blocker bounds yield the global lower bounds.
\item A 55-set construction improves the composition split from 28/28 to 28/27.
\item Lean 4 formalizes both eventual exact saturation numbers.
\end{highlights}

\begin{keywords}
saturated Sperner family \sep Boolean lattice \sep blocker \sep
extremal set theory \sep SAT certificate \sep Lean 4
\end{keywords}

\maketitle

\section{Introduction}

Let $X$ be a finite set.  A family $\family\subseteq\powerset(X)$ is
$k$-Sperner when it contains no strict chain of length $k+1$, and it is
saturated when adjoining any subset outside $\family$ creates such a chain.
The corresponding minimum size, denoted $\sat(n,k)$ when $|X|=n$, asks for a
family that is simultaneously chain-free and maximal with respect to this
property.

Thus, the number 55 counts selected subsets, not points of the ground set.
The difficulty is maximality: every one of the many omitted subsets must
become part of a forbidden eight-term chain as soon as it is adjoined.

Gerbner et al.~\cite{gerbner2013} initiated the systematic study of these
minimum saturated families and proved, among other structural results, that
the value stabilizes when the ground set is sufficiently large for fixed $k$.
The eventual value is denoted $\sat(k)$.  Morrison, Noel, and Scott
\cite{mns2014} developed a layered construction from saturated antichains and
a composition operation that turns an efficient construction for one value of
$k$ into upper bounds for larger values.

Martin and Veldt~\cite{martinveldt2025} used this framework to construct a
56-member saturated 7-Sperner family.  Their construction improved the
composition factor available at $k=7$ and led to the upper bound
$2^{s+1}28^j$ at parameter $5j+2+s$.  They explicitly asked whether
$\sat(7)=56$ or a smaller construction exists.  The same paper records that a
smaller value appearing in early preprint versions was erroneous.  A separate
machine-learning search later produced a 108-member saturated 8-Sperner
family~\cite{patternboost}; that result changes the leading asymptotic bound
but does not answer the question at $k=7$.

There are four main results.  First, \cref{thm:sat6} proves
$\sat(6)=30$, matching the upper construction of Morrison--Noel--Scott.  The
new ingredient is the exact finite-clutter theorem $m(2,3)=9$.  Second,
\cref{thm:sat7-lower} proves $\sat(7)\geq47$ from the stronger local theorem
$m(2,4)=12$.  Third, \cref{sec:exact-seven} excludes every 54-member profile.
Fourth, \cref{sec:construction} gives a 55-member construction.  Together,
these results prove $\sat(7)=55$.

The upper-bound construction has two short, inspectable certificates.  For
each layer, the large core sets are precisely the inclusion-maximal sets that
avoid all small core members.  A general lemma proves saturation, and an
explicit predecessor list proves layering.  A direct finite-ground verifier
and Lean 4 check its concrete eleven-point realization.  The Lean development
also formalizes the arbitrary-ground lower-bound chain and exports the two
eventual exact statements.  Separately, a DRAT-certified encoding proves that
56 is optimal inside the seven-core version of the same template class.  That
class result is not a global lower bound.

A final pre-submission search on 29 July 2026 found no public equivalent of
either exact value in Crossref, arXiv, OpenAlex, Google Scholar, the indexed
citation graph, or the directly relevant full texts.  This bounded negative
search does not establish novelty or priority.

\section{Preliminaries}
\label{sec:preliminaries}

We write $A\subsetneq B$ for strict inclusion.  A \emph{strict chain of
length $r$} is a sequence $A_0\subsetneq\cdots\subsetneq A_{r-1}$.

\begin{definition}
A family $\family\subseteq\powerset(X)$ is \emph{$k$-Sperner} if it has no
strict chain of length $k+1$.  It is \emph{saturated $k$-Sperner} if it is
$k$-Sperner and, for every $S\in\powerset(X)\setminus\family$, the family
$\family\cup\{S\}$ contains a strict chain of length $k+1$.
\end{definition}

A saturated 1-Sperner family is exactly a maximal antichain.  The following
notion organizes several such antichains.

\begin{definition}
Pairwise disjoint families $\mathcal A_0,\ldots,\mathcal A_{k-1}$ are
\emph{layered} if, for every $i>0$ and every $B\in\mathcal A_i$, some
$A\in\mathcal A_{i-1}$ satisfies $A\subsetneq B$.
\end{definition}

We use the following theorem of Morrison--Noel--Scott, in the formulation used
by Martin--Veldt.

\begin{theorem}[Layered saturated antichains
  \cite{mns2014,martinveldt2025}]
\label{thm:layered}
If $\mathcal A_0,\ldots,\mathcal A_{k-1}$ are pairwise disjoint, layered
saturated antichains, then their union is a saturated $k$-Sperner family.
\end{theorem}

Our layers use a common block.  Fix a finite core $C$ and a nonempty set $H$
disjoint from $C$.  A \emph{small template} is a subset $S\subseteq C$; a
\emph{large template} has the form $L\cup H$ for $L\subseteq C$.  A core set
$L$ is \emph{$\smallcores$-free} if it contains no member of
$\smallcores\subseteq\powerset(C)$.

Following Morrison--Noel--Scott, an \emph{atom} of a family is an
inclusion-maximal set $H$ such that every family member either contains $H$
or is disjoint from $H$.  An atom is \emph{homogeneous} when $|H|>2$.

\begin{lemma}[Maximal-small-free certificate]
\label{lem:maximal-free}
Let $\smallcores\subseteq\powerset(C)$ be an antichain, and let
$\largecores$ be the collection of all inclusion-maximal
$\smallcores$-free subsets of $C$.  Then
\[
  \mathcal A(\smallcores,\largecores)
  =\smallcores\ \cup\ \{L\cup H:L\in\largecores\}
\]
is a saturated antichain in $\powerset(C\cup H)$.
\end{lemma}

\begin{proof}
The small templates are mutually incomparable.  Distinct maximal
$\smallcores$-free sets are mutually incomparable, so the large templates are
also mutually incomparable.  A small template $S$ cannot be contained in
$L\cup H$ for $L\in\largecores$, because this would give $S\subseteq L$ and
contradict that $L$ is $\smallcores$-free.  A large template cannot be
contained in a small one because it contains $H$.  Thus $\mathcal A$ is an
antichain.

Now let $T\subseteq C\cup H$ be outside $\mathcal A$, and put $R=T\cap C$.
If $R$ contains some $S\in\smallcores$, then $S\subsetneq T$.  Otherwise $R$
is $\smallcores$-free.  Since $C$ is finite, $R$ is contained in an
inclusion-maximal $\smallcores$-free set $L\in\largecores$, and hence
$T\subseteq L\cup H$.  This containment is strict because $T$ is outside
$\mathcal A$.  Therefore every external set is strictly comparable with a
member of $\mathcal A$, proving saturation.
\end{proof}

\section{Blockers and stable lower bounds}
\label{sec:blockers}

Let $U$ be finite.  A \emph{transversal} of
$\mathcal H\subseteq\powerset(U)$ meets every member of $\mathcal H$, and the
\emph{blocker} $\blocker(\mathcal H)$ is the family of its inclusion-minimal
transversals.  A \emph{clutter} is a family in which no two distinct members
contain one another.  For positive integers $p,q$, define
\[
  m(p,q)=\min\bigl\{|\mathcal S|+|\mathcal C|:
  \mathcal S=\blocker(\mathcal C),\
  \mathcal C=\blocker(\mathcal S),\
  \min_{S\in\mathcal S}|S|\geq p,\
  \min_{C\in\mathcal C}|C|\geq q\bigr\},
\]
where the minimum ranges over all finite ground sets and admissible clutters.

\begin{lemma}[Blocker involution and private witnesses]
\label{lem:blocker-basic}
For every finite clutter $\mathcal H$,
$\blocker(\blocker(\mathcal H))=\mathcal H$.  Moreover, if
$T\in\blocker(\mathcal H)$ and $x\in T$, then some
$E_x\in\mathcal H$ satisfies $T\cap E_x=\{x\}$.  Consequently,
$|T|\leq|\mathcal H|$.
\end{lemma}

\begin{proof}
Minimality of $T$ implies that $T\setminus\{x\}$ misses some $E_x$, while
$T$ meets $E_x$.  Thus $T\cap E_x=\{x\}$.  Distinct points have distinct
private witnesses, proving the cardinality bound.

Every $E\in\mathcal H$ is a transversal of $\blocker(\mathcal H)$.  If a
proper subset $D\subsetneq E$ were also a transversal, then $D$ would contain
no member of the clutter $\mathcal H$.  Hence $U\setminus D$ would meet every
member of $\mathcal H$ and would contain a minimal transversal disjoint from
$D$, a contradiction.  Thus
$\mathcal H\subseteq\blocker(\blocker(\mathcal H))$.  Conversely, every
member of the latter blocker contains a member of $\mathcal H$, and
minimality forces equality.
\end{proof}

The next lemma isolates the interface between canonical layers and blockers.

\begin{lemma}[Canonical blocker interface]
\label{lem:canonical-blocker}
Let $\mathcal A_i$ be an internal canonical layer of a minimum saturated
$k$-Sperner family with homogeneous atom $H$, and put $U=X\setminus H$.
Let $\mathcal S_i$ be the small members of $\mathcal A_i$, and write each
large member as $K\cup H$ with $K\subseteq U$.  Define
\[
  \mathcal C_i=\{U\setminus K:K\cup H\in\mathcal A_i\}.
\]
Then
\[
  \mathcal S_i=\blocker(\mathcal C_i),\qquad
  \mathcal C_i=\blocker(\mathcal S_i).
\]
Every member of $\mathcal S_i$ has size at least $i$, and every member of
$\mathcal C_i$ has size at least $k-1-i$.
\end{lemma}

\begin{proof}
The size statements are Proposition 14 of Martin--Veldt
\cite{martinveldt2025}.  Choose a nonempty proper subset $Q\subset H$.  For
each $R\subseteq U$, the mixed set $R\cup Q$ is outside the layer.  Saturation
of $\mathcal A_i$ gives either $S\subseteq R$ for some
$S\in\mathcal S_i$, or $R\subseteq K$ for some large trace $K$.  The
antichain property says that every $S$ meets every $U\setminus K$.

If $T$ is a transversal of $\mathcal C_i$, it is contained in no large trace,
so the cover condition gives some $S\subseteq T$.  Minimal transversals are
therefore precisely the members of $\mathcal S_i$.  Applying the same argument
to $U\setminus T$ gives the reverse blocker equality.  Alternatively, it
follows from \cref{lem:blocker-basic}, since both sides are clutters.
\end{proof}

\begin{theorem}
\label{thm:m23}
The exact local blocker parameter is $m(2,3)=9$.
\end{theorem}

\begin{proof}
Consider an admissible pair and write
$a=|\mathcal S|$, $b=|\mathcal C|$.  The private-witness bound applied in both
directions gives $a\geq3$ and $b\geq2$.

If $b=2$, the two members of $\mathcal C$ are disjoint, because their
intersection would give a singleton blocker.  Choosing one point from each
gives at least $3^2=9$ members of $\mathcal S$.

Suppose $b=3$, with rows $C_1,C_2,C_3$.  Their triple intersection is empty.
For $\{i,j,k\}=\{1,2,3\}$, put
\[
  A_i=C_i\setminus(C_j\cup C_k),\qquad
  P_i=(C_j\cap C_k)\setminus C_i,
\]
and write $\alpha_i=|A_i|$, $\pi_i=|P_i|$.  The two-point minimal
transversals use one point from $A_i$ and one from $P_i$, or one point from
each of two distinct $P$-classes.  The three-point minimal transversals use
one point from each $A_i$.  The private-witness part of
\cref{lem:blocker-basic} excludes every other type.  Hence
\[
  a=\alpha_1\alpha_2\alpha_3+
  \sum_{i=1}^3\alpha_i\pi_i+
  \pi_1\pi_2+\pi_1\pi_3+\pi_2\pi_3.
  \tag{1}
\]
The row inequalities are
$\alpha_i+\pi_j+\pi_k\geq3$.  If some $\pi_i\geq2$, the terms in (1)
containing $\pi_i$ give at least six.  Equality would force
$\pi_i=2$, the corresponding row sum to be three, and every remaining term
to vanish.  The other two row inequalities then force all three
$\alpha$-classes to be nonempty, a contradiction.  Thus $a\geq7$.  If every
$\pi_i\leq1$, monotonicity in the $\alpha_i$ reduces the calculation, up to
symmetry, to
\[
\begin{array}{c|c|c}
(\pi_1,\pi_2,\pi_3)&
\min(\alpha_1,\alpha_2,\alpha_3)&\min a\\
\hline
(0,0,0)&(3,3,3)&27\\
(0,0,1)&(2,2,3)&15\\
(0,1,1)&(1,2,2)&9\\
(1,1,1)&(1,1,1)&7
\end{array}
\]
so again $a\geq7$.

It remains to consider $b\geq4$.  If $a=3$, the three members of
$\mathcal S$ are pairwise disjoint: an intersection point and one point from
the remaining row would otherwise form a transversal of size at most two.
Choosing one point from each row gives $b\geq2^3=8$.  If $a=4$, form the
intersection graph on the four rows of $\mathcal S$.  Two vertex-disjoint
edges would again give a transversal of size at most two, so the matching
number is at most one.  With two isolated rows, blocker factorization gives at
least $2\cdot2\cdot2=8$ transversals.  With one isolated row, the connected
three-row component has at least three blockers and factorization gives at
least $2\cdot3=6$.  With no isolated row, the graph is a three-leaf star.
The leaves are pairwise disjoint; choosing one point from each gives at least
\[
  n_1n_2n_3-(n_1-1)(n_2-1)(n_3-1)\geq7
\]
minimal transversals, where $n_i\geq2$ are the leaf sizes.  Hence $b\geq6$.
If $a\geq5$, the desired total is immediate.  These cases prove
$a+b\geq9$.

For equality, take distinct points $x_1,x_2,y_1,y_2,w,z$ and set
\[
  \mathcal S=\{\{x_i,y_j\}:i,j\in\{1,2\}\}\cup\{\{w,z\}\},
\]
\[
  \mathcal C=\{\{x_1,x_2,w\},\{x_1,x_2,z\},
  \{y_1,y_2,w\},\{y_1,y_2,z\}\}.
\]
A minimal transversal of $\mathcal C$ either chooses one $x_i$ and one
$y_j$, or equals $\{w,z\}$.  Thus
$\blocker(\mathcal C)=\mathcal S$, and blocker involution gives
$\blocker(\mathcal S)=\mathcal C$.  The total size is nine.
\end{proof}

\begin{theorem}
\label{thm:sat6}
The eventual saturated 6-Sperner number is
\[
  \sat(6)=30.
\]
\end{theorem}

\begin{proof}
Morrison--Noel--Scott constructed a 30-member saturated 6-Sperner family on
every ground set of size at least eight
\cite[Proposition~20]{mns2014}, so $\sat(6)\leq30$.

Choose $n$ in the stable range with $n\geq8$ and $n>2^{30}$, and choose an
endpoint-normal minimum family $\family$ on an $n$-element set.  Since
$|\family|\leq30$, the homogeneous-atom theorem applies.  The canonical
decomposition gives six nonempty saturated-antichain layers
$\mathcal A_0,\ldots,\mathcal A_5$.  The endpoint result and
\cite[Lemma~15]{martinveldt2025} give
\[
  |\mathcal A_0|=|\mathcal A_5|=1,\qquad
  |\mathcal A_1|,|\mathcal A_4|\geq5.
\]
By \cref{lem:canonical-blocker}, the second layer is an admissible
$(2,3)$ blocker pair, and the third becomes one after exchanging the two
sides.  Thus \cref{thm:m23} gives
$|\mathcal A_2|,|\mathcal A_3|\geq9$.  Therefore
\[
  |\family|\geq1+5+9+9+5+1=30.
\]
The upper and lower bounds agree.
\end{proof}

\section{A global lower bound at seven layers}
\label{sec:sat7-lower}

The technical blocker estimate needed at seven layers is proved in
\cref{app:m24}.

\begin{theorem}
\label{thm:m24}
The exact local blocker parameter is $m(2,4)=12$.
\end{theorem}

\begin{proof}
By \cref{prop:m24-product}, every admissible pair satisfies
$ab\geq32$, where $a$ and $b$ are its side cardinalities.  Hence
\[
  (a+b)^2=(a-b)^2+4ab\geq128>11^2,
\]
so $a+b\geq12$.

For equality, let $\mathcal S$ be the eight edges of two vertex-disjoint
four-cycles.  The blocker of each cycle consists of its two bipartition
classes, so $\blocker(\mathcal S)$ consists of the four unions of one class
from each cycle.  Conversely, any transversal of these four unions contains
an edge of one of the cycles.  Thus the two families are mutual blockers,
their member-size minima are two and four, and their total cardinality is
$8+4=12$.
\end{proof}

\begin{lemma}
\label{lem:middle-nine}
In the canonical decomposition used for a minimum saturated 7-Sperner
family, the middle layer satisfies $|\mathcal A_3|\geq9$.
\end{lemma}

\begin{proof}
Use the notation of \cref{lem:canonical-blocker}.  Every member of
$\mathcal S_3$ and $\mathcal C_3$ has size at least three.  Choose
$R\subseteq U$ uniformly.  The canonical cover says that at least one event
$S\subseteq R$ or $R\cap C=\varnothing$ occurs.  Each event has probability
at most $2^{-3}$.  The union bound gives
$|\mathcal A_3|\geq8$.

If equality held, every event would have probability $1/8$ and all events
would be pairwise disjoint.  The choice $R=U$ belongs to every small-side
event, while $R=\varnothing$ belongs to every large-complement event.  Both
sides are nonempty, so equality would force one event on each side, contrary
to a total of eight.  Hence $|\mathcal A_3|\geq9$.
\end{proof}

\begin{theorem}
\label{thm:sat7-lower}
The eventual saturated 7-Sperner number satisfies
\[
  \sat(7)\geq47.
\]
\end{theorem}

\begin{proof}
Choose a stable ground set with $n>2^{55}$.  The construction in
\cref{thm:size55} ensures that a minimum family has at most 55 members, so a
homogeneous atom and the canonical decomposition are available.  The endpoint
and adjacent-layer bounds give
\[
  (|\mathcal A_0|,|\mathcal A_1|,|\mathcal A_5|,
  |\mathcal A_6|)\geq(1,6,6,1).
\]
The canonical blocker interface and \cref{thm:m24} give
$|\mathcal A_2|,|\mathcal A_4|\geq12$, while
\cref{lem:middle-nine} gives $|\mathcal A_3|\geq9$.  Summing the seven layers
gives
\[
  \sat(7)\geq1+6+12+9+12+6+1=47.
\]
\end{proof}

\section{The 55-member construction}
\label{sec:construction}

Let $C=[8]=\{1,\ldots,8\}$, let $H$ be disjoint from $C$, and assume
$|H|>2$.  We abbreviate a core subset by concatenating its elements; for
example, $145$ denotes $\{1,4,5\}$ and $1346H$ denotes
$\{1,3,4,6\}\cup H$.  The symbol $-$ denotes an empty list.

For each row of \cref{tab:layers}, let $\smallcores_i$ be the listed small
cores and let $\largecores_i$ be the listed large cores.  Define
\[
  \mathcal A_i
   =\smallcores_i\cup\{L\cup H:L\in\largecores_i\},
  \qquad 0\leq i\leq6.
\]

\begin{table}[htbp]
\centering
\small
\setlength{\tabcolsep}{4pt}
\renewcommand{\arraystretch}{1.18}
\begin{tabular}{c p{0.35\textwidth} p{0.48\textwidth}}
\toprule
$i$ & $\smallcores_i$ & $\largecores_i$ \\
\midrule
0 & $\varnothing$ & -- \\
1 & $1$, $2$, $3$, $6$, $7$ & $458$ \\
2 & $24$, $15$, $35$, $26$, $47$, $67$, $18$, $38$
  & $1346$, $1237$, $4568$, $2578$ \\
3 & $124$, $234$, $145$, $345$, $156$, $356$, $267$, $178$, $378$
  & $12357$, $13467$, $12358$, $12368$, $13468$, $24568$, $24578$, $45678$ \\
4 & $2345$, $1247$, $3568$, $1678$
  & $123567$, $134567$, $123468$, $124568$, $123578$, $134578$, $234678$, $245678$ \\
5 & $23457$
  & $1234568$, $1234678$, $1235678$, $1245678$, $1345678$ \\
6 & -- & $12345678$ \\
\bottomrule
\end{tabular}
\caption{Core description of the seven layers.  Every large core in the
right column represents its union with $H$.}
\label{tab:layers}
\end{table}

\begin{proposition}
\label{prop:layer-certificate}
For every $i\in\{0,\ldots,6\}$, the family $\mathcal A_i$ is a saturated
antichain.
\end{proposition}

\begin{proof}
Within each nonempty small column of \cref{tab:layers}, all sets have the same
cardinality, so $\smallcores_i$ is an antichain.  For each row, the large-core
column is exactly the collection of inclusion-maximal core subsets that avoid
all members of the small-core column.  This statement is a finite certificate:
adding any missing core element to a listed large core creates one of the
listed small cores, while every small-free core set is contained in one of the
listed large cores.  The latter coverage implication has only $2^8$ core
masks per row; the supplementary verifier checks both implications for all
seven rows.  The endpoint rows say respectively that no set avoids
$\varnothing$ and that $C$ is the unique maximal set avoiding an empty small
family.  The result now follows from \cref{lem:maximal-free}.
\end{proof}

For completeness, \cref{app:layering} records one predecessor for every member
above layer zero.  Each displayed containment is strict, so the seven layers
are layered.  The 55 templates in \cref{tab:layers} are pairwise distinct
across the rows.

\begin{theorem}
\label{thm:size55}
For every finite $H$ disjoint from $[8]$ with $|H|>2$, the family
\[
  \family=\bigcup_{i=0}^{6}\mathcal A_i
\]
is a saturated 7-Sperner family of cardinality 55.
\end{theorem}

\begin{proof}
By \cref{prop:layer-certificate}, all seven layers are saturated antichains.
The predecessor certificate in \cref{app:layering} proves layering, and direct
inspection of \cref{tab:layers} proves pairwise disjointness.  Thus
\cref{thm:layered} shows that their union is saturated 7-Sperner.  The layer
sizes are
\[
  1,\ 6,\ 12,\ 17,\ 12,\ 6,\ 1,
\]
whose sum is 55.
\end{proof}

\begin{proposition}
\label{prop:endpoints-atom}
The family $\family$ contains both $\varnothing$ and $C\cup H$, and $H$ is a
homogeneous atom of $\family$.  Relative to this atom, $\family$ has 28 small
members and 27 large members.
\end{proposition}

\begin{proof}
The two endpoints are the unique members of $\mathcal A_0$ and
$\mathcal A_6$.  Every displayed template either contains all of $H$ or is
disjoint from $H$.  To see maximality, note that every core point occurs in a
small member: points $1,2,3,6,7$ occur in singleton members of
$\mathcal A_1$, while $4,5,8$ occur respectively in $24$, $15$, and $18$.
Thus, for each $c\in C$, some member of $\family$ meets $H\cup\{c\}$ in a
nonempty proper subset.  No strict superset of $H$ can therefore have the
atom property.  Since $|H|>2$, the atom is homogeneous.  Finally, counting
the two columns of \cref{tab:layers} gives
$1+5+8+9+4+1=28$ small templates and
$1+4+8+8+5+1=27$ large templates.
\end{proof}

\par\medskip

\begin{corollary}
\label{cor:main-bound}
For every $n\geq11$, one has $\sat(n,7)\leq55$.  Consequently,
$\sat(7)\leq55$.
\end{corollary}

\begin{proof}
Choose $|H|=n-8$.  Then $|H|>2$, so \cref{thm:size55} applies on an
$n$-element ground set.  The eventual inequality follows from the definition
of $\sat(7)$.
\end{proof}

\section{Excluding 54 members}
\label{sec:exact-seven}

We now refine the lower bound.  Let
$\mathcal A_0,\ldots,\mathcal A_6$ be the canonical layers of a minimum
family on a stable ground set with a homogeneous atom.  For an internal layer,
write $\mathcal S_i$ for its small traces and $\mathcal C_i$ for the
complements of its large traces, as in \cref{lem:canonical-blocker}.

\begin{proposition}[Canonical profile reduction]
\label{prop:canonical-profiles}
No canonical family has total size 52 or 53.  If a canonical family has total
size 54, its layer profile is one of
\[
  P_{\mathrm F}=(1,6,13,14,13,6,1),\qquad
  P_{16}=(1,6,12,16,12,6,1).
\]
In the first profile the middle blocker pair is the Fano line clutter.  In the
second profile the middle blocker pair has total cardinality 16.
\end{proposition}

\begin{proof}
The endpoint, adjacent-layer, and blocker bounds give
\[
  (|\mathcal A_0|,\ldots,|\mathcal A_6|)
  \geq(1,6,12,14,12,6,1).
  \tag{2}
\]
The middle bound follows by applying private witnesses, residual blockers,
and the three-row incidence formula to every blocker split of total at most
13.  When every point has degree at most two, the active points form an actual
loop-multigraph on the fixed row labels.  Inclusion-minimality forces every
positive support to meet a minimum-degree row, so the row-size bounds leave a
finite exhaustive support problem.  Independent implementations agree on all
saved classifications.

Equality in (2) forces the middle pair to be the Fano plane.  Its lower or
upper adjacent 12-member equality interface is impossible, so total 52 is
excluded.  With one slack member, either the same Fano exclusion applies or
the middle pair has total 15.  The splits $(5,10)$, $(6,9)$, $(7,8)$ and their
duals are excluded by the same arbitrary-finite degree and residual
reductions, so total 53 is excluded.

For total 54, write the five internal slacks over (2) as nonnegative integers
with sum two.  A 14-member middle pair is Fano and forces both adjacent layers
above equality.  A 15-member middle pair is impossible.  If the middle layer
has more than 15 members, it has exactly 16 and every other layer attains
equality.  These alternatives give exactly the two displayed profiles.
\end{proof}

For the exact exclusion, put
\[
\begin{aligned}
  \mathcal F&=\mathcal S_2,&\mathcal G&=\mathcal C_2,\\
  \mathcal H&=\mathcal S_3,&\mathcal K&=\mathcal C_3,\\
  \mathcal P&=\mathcal S_4,&\mathcal Q&=\mathcal C_4.
\end{aligned}
\]
The adjacent interfaces say that every member of $\mathcal G$ strictly
contains a member of $\mathcal K$, with
$\mathcal F\cap\mathcal H=\varnothing$, and dually that every member of
$\mathcal P$ strictly contains a member of $\mathcal H$, with
$\mathcal K\cap\mathcal Q=\varnothing$.

\begin{proposition}[Fifty-four-member exclusion]
\label{prop:exclude54}
Neither profile in \cref{prop:canonical-profiles} is realizable on any finite
trace ground set.
\end{proposition}

\begin{proof}
For $P_{\mathrm F}$, each Fano line strictly contains a distinct pair in
$\mathcal F$.  Up to blocker duality, the adjacent split is $(7,6)$,
$(8,5)$, or $(9,4)$.  In the first split, $\mathcal F$ chooses one of three
pairs on each Fano line.  All $3^7=2187$ choices are checked, and none has six
minimal vertex covers of size at least four.  In the second split, write
$\mathcal F=E(Q)\cup\{R\}$.  The trace $R\cap X$ is independent in $Q$, while
the outside points of $R$ form one incidence class whose exact multiplicity
is at most five by the private-witness bound.  This reduces every finite
instance to 249576 triples $(Q,R\cap X,t)$.  No admissible triple has five
blockers of minimum size four.  In the last split, the four rows of
$\mathcal G$ have size at least four and intersect pairwise.  A residual trace
argument gives at least ten blockers, contradicting $|\mathcal F|=9$.
Thus the Fano profile and its upper dual are impossible.

For $P_{16}$, a uniformly random subset gives an exhaustive disjoint cover by
the events $D\subseteq R$ and $E\subseteq U\setminus R$ for the two middle
blocker sides.  If all 16 rows had size at least four, the union bound would be
an equality, but two events on one side overlap with positive probability.
Hence some middle row has size three.

Up to exchanging blocker sides, the middle split is $(5,11)$, $(6,10)$,
$(7,9)$, or $(8,8)$.  The following table records the exhaustive reductions.
All multiplicities count actual points, so parallel supports and repeated
incidence patterns remain distinct.
\begin{table}[htbp]
\centering
\small
\setlength{\tabcolsep}{5pt}
\renewcommand{\arraystretch}{1.16}
\begin{tabular}{c p{0.41\textwidth} p{0.38\textwidth}}
\toprule
Split & Arbitrary-finite reduction & Exhaustive residue \\
\midrule
$(5,11)$ & A degree-three point gives nine blockers containing it and two
residual blockers avoiding it. & Equality would make the residual the blocker
of two disjoint rows of size at least three, producing at least nine residual
rows from only five. \\
$(6,10)$ & Degrees at least four give eleven blockers.  Degree three leaves
three exact three-row kernels; degree at most two gives an actual
loop-multigraph. & The nonanalytic kernel has 95 multiplicity vectors and zero
compatible extensions.  Every low-degree core has at least 15 blockers. \\
$(8,8)$ & Degree reductions leave maximum degree three.  A three-point row
then leaves a five-row kernel with rank-at-most-three supports. & There are
1890 multiplicity vectors in 24 classes and 2514 possible extensions; none is
valid. \\
$(7,9)$ & Maximum degrees are three and four.  A three-point row on the
nine-row side leaves a four-row kernel; if it occurs only on the seven-row
side, two three-row residuals must be coupled. & The first branch has 136
vectors, 138161 extensions, and no valid case.  All 90898 couplings in the
second branch have at least four pair blockers, contradicting the bound three.
\\
\bottomrule
\end{tabular}
\caption{Exhaustive middle-split reductions for the 16-member profile.}
\label{tab:exact-splits}
\end{table}

Each finite domain in \cref{tab:exact-splits} follows from row count, private
witnesses, point-degree bounds, and support minimization; no finite bound on
the original trace ground set is assumed.  C++ enumerators and independent
Python implementations reproduce the saved outputs byte for byte.  Thus all
four splits and their blocker duals are impossible.  The forced three-point
row has nowhere to occur, excluding $P_{16}$.
\end{proof}

\begin{theorem}
\label{thm:sat7}
The eventual saturated 7-Sperner number is
\[
  \sat(7)=55.
\]
\end{theorem}

\begin{proof}
Choose a minimum family in the stable range.  The 55-member construction
ensures that the homogeneous-atom reduction applies.  The total-52 and
total-53 exclusions give $\sat(7)\geq54$, and
\cref{prop:exclude54} excludes equality.  Hence $\sat(7)\geq55$.
The reverse inequality is \cref{cor:main-bound}.
\end{proof}

\section{Composition consequences}
\label{sec:composition}

By \cref{prop:endpoints-atom}, the family in \cref{thm:size55} satisfies the
endpoint and homogeneous-atom hypotheses of the composition lemma of
Morrison--Noel--Scott~\cite[Lemma~18]{mns2014}.  That lemma combines two saturated
systems on disjoint ground sets by taking small--small and large--large
pairs.  If the factors have small/large sizes $(a_1,b_1)$ and $(a_2,b_2)$,
the composition has sizes $(a_1a_2,b_1b_2)$ and saturation parameter
$k_1+k_2-2$; it also preserves the endpoint and homogeneous-atom hypotheses.

Iterating the 55-member factor $j\geq1$ times therefore gives parameter
$5j+2$ and small/large sizes $(28^j,27^j)$.  For the remainder factor, take
an $s$-element set $Y$ and another homogeneous atom $K$, and use
\[
  \powerset(Y)\ \cup\ \{S\cup K:S\subseteq Y\}.
\]
This is the standard saturated $(s+2)$-Sperner system: its small and large
parts both have size $2^s$, and it contains both endpoints.  Composing it with
the iterated factor changes the parameter by
$(s+2)-2=s$ and gives small/large sizes
$2^s28^j$ and $2^s27^j$, respectively.  We use $0\leq s\leq4$ to select the
usual five residue classes.

\begin{corollary}
\label{cor:discrete-bound}
For $j\geq1$ and $0\leq s\leq4$,
\[
  \sat(5j+2+s)\leq 2^s\bigl(28^j+27^j\bigr).
\]
\end{corollary}

For comparison, the symmetric 56-member factor gives
$2^{s+1}28^j$.  The inequality
\[
  2^s(28^j+27^j)<2^{s+1}28^j
\]
is strict for every $j\geq1$.  The leading exponential rate is unchanged,
because the larger part still has size 28.

\section{Seven-core class optimality}
\label{sec:class-optimality}

The 55-member construction uses an eight-element core.  We also tested whether
the analogous seven-core template class can contain 55 members.  The next
result is computer-assisted and has a deliberately narrower scope than
\cref{thm:size55}.

\begin{proposition}[Certificate-backed class optimum]
\label{prop:class-optimum}
Let $|C|=7$ and $|H|>2$.  Among all ordered seven-tuples of pairwise disjoint,
layered saturated antichains whose members have the form $A$ or $A\cup H$ with
$A\subseteq C$, the minimum total cardinality is 56.
\end{proposition}

\begin{proof}[Computer-assisted proof]
For layer $i$, atom-use bit $b$, and core mask $m$, introduce a Boolean
variable $x_{i,b,m}$ specifying whether the associated template is selected.
The CNF contains: exclusions for comparable pairs within a layer; one coverage
clause for each of the three homogeneous-block signatures and each core mask;
pairwise exclusions between layers; predecessor implications for layering;
and a sequential counter~\cite{sinz2005} restricting the total number of
selected templates to at most 55.

The three signatures record whether an arbitrary test set meets $H$ in
$\varnothing$, a nonempty proper subset, or all of $H$.  Inclusion against a
template depends only on this signature and the core mask.  Hence every
mathematical candidate of size at most 55 extends to a satisfying assignment
of the CNF.  Conversely, the selected variables of a satisfying assignment
decode to the stated template conditions; independent semantic checks are
included as controls.

The deterministic instance contains 100,297 variables and 252,427 clauses.
Kissat 4.0.4 returned UNSAT and emitted a 253,731,315-byte DRAT proof.
DRAT-trim~\cite{wetzler2014} independently checked the proof and returned
\texttt{s VERIFIED}.  Therefore no seven-core candidate of size at most 55
exists.  The 56-member construction of Martin--Veldt supplies the matching
upper bound.
\end{proof}

The encoding imposes neither complement symmetry nor prescribed layer sizes,
and it does not require the Fano-plane pattern.  Nevertheless,
\cref{prop:class-optimum} does not cover larger cores, non-template families,
or systems without a common block, and therefore gives no global lower bound
on $\sat(7)$.

\section{Independent verification and formalization}
\label{sec:verification}

The discovery computations, independently replayed certificates, and Lean
formalization have different logical roles.  The proof of
\cref{thm:size55} is the structural argument from
\cref{lem:maximal-free} and the displayed certificates; it does not trust the
SAT solver used to discover the family.  Four redundant exact checks were
also performed.

\begin{enumerate}[leftmargin=*]
  \item A symbolic verifier checks all $3\cdot2^8=768$ core/block signatures
        for every layer.
  \item Expanding $H$ to three points gives an eleven-element ground set; a
        second verifier checks each layer on the full powerset.
  \item An integer-bitmask dynamic program checks that the 55-set union has no
        chain of length eight and checks a witness after adjoining each of the
        other 1993 ground subsets.
  \item A separately implemented set/DAG verifier repeats the union-level
        check without sharing the bitmask dynamic program.
\end{enumerate}

All four checks accept the same semantic family; its SHA-256 identifier is
recorded in the supplementary computational certificate.
As a negative control, they reject the withdrawn 54-member preprint family at
its third layer, with core set $\{2,3,6\}$ uncovered.

The concrete eleven-element realization was formalized in Lean
4~\cite{lean4} using Mathlib~\cite{mathlib2020}.  Lean proves that the explicit
family has cardinality 55, is 7-Sperner, and is saturated.  Saturation is
checked through a frozen table covering all $2^{11}=2048$ masks: 55 member
rows and 1993 external rows containing a strict chain of length eight after
insertion.  The table is divided into 32 kernel-reduced blocks.  This theorem
covers the concrete case $H\cong[3]$; the arbitrary-$H$ construction in
\cref{thm:size55} is proved independently by the maximal-small-free argument.

The final Lean development goes beyond this concrete certificate and
formalizes both eventual exact values.  Its central predicate is not a
numeric function evaluated by Lean.  Writing
$\operatorname{Sat}_k(\mathcal F)$ for the assertion that $\mathcal F$ is a
saturated $k$-Sperner family, the proposition
\texttt{IsStableSaturationNumber k s} expands to
\[
  \exists N\ \forall n\geq N,\quad
  \left(\exists\mathcal F\subseteq\powerset([n]),
    \operatorname{Sat}_k(\mathcal F)\ \land\ |\mathcal F|=s\right)
  \land
  \left(\forall\mathcal F\subseteq\powerset([n]),
    \operatorname{Sat}_k(\mathcal F)\ \Longrightarrow\ s\leq|\mathcal F|\right).
\]
Thus the proposition contains, after one threshold, both an upper-bound
witness and a universal lower bound on every labelled ground
\texttt{Fin n}.  In their respective stable-exact namespaces, the two
exported declarations are
\begin{quote}\small\ttfamily
sat\_six\_eq\_thirty : IsStableSaturationNumber 6 30\\
sat\_seven\_eq\_fifty\_five : IsStableSaturationNumber 7 55.
\end{quote}

For six layers, the formal proof combines a parameterized 30-member
construction with the arbitrary-finite theorem $m(2,3)=9$, homogeneous-atom
existence, canonical decomposition, and the resulting layer sum.  For seven
layers, it combines the parameterized 55-member construction with the
canonical profile reduction, the Fano-adjacent exclusion, and every middle
split of total 16.  In particular, the Fano, $(5,11)$, $(6,10)$, $(7,9)$,
and $(8,8)$ branches and their blocker duals now occur in the dependency graph
of the exported exact theorem.  The earlier local theorem
\texttt{no\_kernel\_completion\_of\_eight} remains the arbitrary-finite
five-row component of the $(8,8)$ branch.

The finite branches in \cref{prop:canonical-profiles,prop:exclude54} retain a
separate manifest-bound trust chain.  The G4.10, G4.11, and G4.12 manifests
bind 8, 23, and 36 proof, source, and output objects, respectively.  Every
saved enumeration was replayed, and independent implementations agreed on
the semantic counts.  In Lean, executable checks are connected to the
mathematical statements by proved soundness and realization lemmas; a raw
search result is never used as an unproved theorem.

The final formalization audit was completed on 27 July 2026.  The full Lake
build reported successful completion of 17,488 jobs, and the independent
entry point \nolinkurl{Problems/P0054/formal/Main.lean} checked both exact
theorems successfully.  For the exported exact statements,
\texttt{\#print axioms} reports only \texttt{propext},
\texttt{Classical.choice}, and \texttt{Quot.sound}.  A source-wide scan of
the Lean files found zero occurrences of \texttt{sorry}, \texttt{admit}, an
\texttt{axiom} declaration, \texttt{unsafe}, \texttt{native\_decide}, or
\texttt{run\_tac}.

We use three evidence labels throughout the project.  \texttt{COMPUTED}
denotes discovery searches and redundant semantic replays that are not, by
themselves, proof dependencies.  \texttt{PROVED} denotes the mathematical
arguments and independently checked certificate reductions in this
manuscript.  \texttt{FORMALIZED} denotes statements accepted by the Lean
kernel, including $m(2,3)=9$, $m(2,4)=12$, the concrete 55-member
certificate, the branch exclusions, and the two eventual exact values.  These
labels concern logical support; none establishes novelty, priority, or
independent peer review.

\section{Discussion and open problems}

The construction and exclusion theorem answer the numerical question posed by
Martin--Veldt: the eventual value is 55.  The construction's asymmetry is
useful under composition.  Although replacing the split $28/28$ by $28/27$
does not change the leading asymptotic exponent, it improves every finite
bound in \cref{cor:discrete-bound}.

The seven-core class result gives a limited structural explanation for why the
published family could not simply be optimized within its original core size:
the complete common-block layered template class on seven core points has
minimum 56.  The 55-member family escapes by using an eighth core point.  This
does not imply that every smaller global construction must use eight core
points or even admit a homogeneous-block description.

The exact six-layer value closes the gap left by the published 30-member
construction.  Its blocker proof is independent of the finite
incidence-pattern search, which serves only as a regression and
counterexample check.

Kernel acceptance complements rather than replaces human reconstruction of
the arbitrary-finite reductions and scrutiny of the certificate interfaces.
The final pre-submission search found no public equivalent in its recorded
scope.  This negative result remains bounded and does not establish novelty
or priority.

Another open direction is the proposed identity
$\sat(k)=\operatorname{A075529}(k)$, where A075529 counts distinct terminal
values of star addition chains.  This conjecture remains \texttt{UNKNOWN}:
the agreement of the first seven values neither proves a general formula nor
supplies a structural correspondence.  We defer its study to future work.

\section*{Data and code availability}

This work is theoretical and does not generate or analyse empirical datasets.
The supporting artifact package contains the Lean 4 formalization, exact
search programs, construction verifiers, certificate generators, CNF
instances, DRAT proofs, verification scripts, manifests, checksums, tests, and
locked dependency metadata.  The package is available from the versioned
GitHub release
\href{https://github.com/lailai0916/saturated-sperner-6-7/releases/tag/v1.0.0}
{v1.0.0} and is archived on Zenodo under DOI
\href{https://doi.org/10.5281/zenodo.21679078}{10.5281/zenodo.21679078}.

\section*{Declaration of generative AI and AI-assisted technologies in the
manuscript preparation process}

During the preparation of this work, the author used OpenAI Codex for
literature-query design, proof and code drafting, verification orchestration,
and language editing.  The author reviewed and edited all generated content
and takes full responsibility for the content of the article.  The
mathematical claims rely on the displayed arguments, independently replayable
certificates, direct verifiers, and Lean kernel acceptance rather than on
model output.

\section*{Funding}

This research did not receive any specific grant from funding agencies in the
public, commercial, or not-for-profit sectors.

\section*{Declaration of competing interest}

The author declares that he has no known competing financial interests or
personal relationships that could have appeared to influence the work
reported in this paper.

\printcredits

\appendix

\section{A product inequality for a local blocker parameter}
\label{app:m24}

\begin{proposition}
\label{prop:m24-product}
Let $\mathcal S$ and $\mathcal C$ be mutual blockers on a finite ground set,
with every member of $\mathcal S$ of size at least two and every member of
$\mathcal C$ of size at least four.  Then
\[
  |\mathcal S||\mathcal C|\geq32.
\]
\end{proposition}

\begin{proof}
Put $a=|\mathcal S|$ and $b=|\mathcal C|$.  The private-witness lemma gives
$a\geq4$ and $b\geq2$.  We prove the following six bounds:
\[
\begin{array}{c|c@{\qquad}c|c}
b&\text{lower bound on }a&a&\text{lower bound on }b\\
\hline
2&16&4&16\\
3&12&5&8\\
4&8&6&6
\end{array}
\tag{2}
\]

For $b=2$, the two members of $\mathcal C$ are disjoint, and choosing one
point from each gives at least $4^2=16$ blockers.  For $b=3$, use the six
incidence classes from the proof of \cref{thm:m23}.  Formula (1) remains
valid, now with
$\alpha_i+\pi_j+\pi_k\geq4$.  If some $\pi_i\geq3$, its terms contribute at
least $3\cdot4=12$.  Otherwise, order
$0\leq\pi_1\leq\pi_2\leq\pi_3\leq2$ and minimize each $\alpha_i$.
The ten possibilities give respectively
\[
  64,40,24,25,16,12,17,13,12,12,
\]
so $a\geq12$.  If $a=4$, every transversal of $\mathcal S$ has size at least
four.  Two intersecting rows, together with one point from each remaining
row, would give a transversal of size at most three.  Hence the four rows are
pairwise disjoint, and choosing one point from each gives $b\geq2^4=16$.

We next prove the $b=4$ bound.  Write
$\mathcal C=\{C_1,C_2,C_3,C_4\}$ and choose a minimum-size member
$E=C_1$, with $n=|E|\geq4$.  For $j=2,3,4$, set
\[
  O_j=E\setminus C_j,\qquad R_j=C_j\setminus E,\qquad
  o_j=|O_j|,\quad q_j=|R_j|.
\]
The clutter property makes all $O_j,R_j$ nonempty, minimality of $|E|$ gives
$q_j\geq o_j$, and the absence of singleton blockers gives
$O_2\cup O_3\cup O_4=E$.  For $x\in E$, let
\[
  \mathcal J_x=\{R_j:x\in O_j\}.
\]
The map $R\mapsto\{x\}\cup R$ is a bijection from
$\blocker(\mathcal J_x)$ to the members of $\blocker(\mathcal C)$ whose
intersection with $E$ is $\{x\}$.  This follows in both directions from
private witnesses.  Consequently, if
$w_x=\min_{j:x\in O_j}o_j$, then
\[
  |\blocker(\mathcal C)|\geq\sum_{x\in E}w_x.
\tag{3}
\]
Assign each $x$ to an index attaining $w_x$, and let $t_j$ be the three
class sizes.  Since $t_j\leq o_j$ and $\sum t_j=n$,
\[
  \sum_{x\in E}w_x=\sum_{j=2}^4t_jo_j
  \geq\sum_{j=2}^4t_j^2
  \geq\left\lceil\frac{n^2}{3}\right\rceil.
\tag{4}
\]
For $n\geq5$, this is at least nine.  Let $n=4$.  The right side of (3) is
at least six.  If it is seven, the nonempty traces
$E,E\setminus O_2,E\setminus O_3,E\setminus O_4$ have a minimal
transversal of size at least two, which gives one further blocker outside all
singleton-intersection classes.  If it is six, equality in (4) forces the
$O_j$ to partition $E$ with sizes $2,1,1$; the four cross-pairs between the
two-point block and the singleton blocks, together with the pair of
singletons, give five further blockers.  Thus
$|\blocker(\mathcal C)|\geq8$ in every case.

For $a=5$, every transversal of $\mathcal S$ has size at least four.  A point
belongs to at most two rows.  The row-intersection graph has matching number
at most one, since two disjoint edges and one point from the fifth row would
give a three-point transversal.  Such a graph is empty, a star, or a triangle,
with isolated vertices allowed.  The empty case gives $2^5$ blockers.  In a
star, fix one center--leaf intersection point and choose one point from each
of the other three pairwise disjoint rows.  In a triangle, fix an intersection
point on one edge and choose one point from the third triangle row and each
isolated row.  Incidence degree at most two supplies a private witness for the
fixed point.  Both nonempty cases give at least $2^3=8$ blockers.

Finally, let $a=6$.  If a point lies in three rows, the other three rows are
pairwise disjoint; fixing that point and choosing one point from each remaining
row gives at least eight blockers.  Otherwise encode each actual ground point
as a loop or edge on the six row labels according as it belongs to one or two
rows.  Parallel points remain distinct.  Each vertex has incidence degree at
least two, and the nonloop support graph has matching number at most two.

If the matching number is zero, loop choices give at least $2^6$ minimal
covers.  If it is one, the nonloop support is a star or a triangle, with
isolated vertices allowed.  In the star case, fix one actual center--leaf
point and choose one incident point at each of the other four vertices.  No
nonloop support joins two of those vertices, so the choices are distinct and
each has a private vertex; this gives at least $2^4$ covers.  In the triangle
case, fix actual points on two adjacent edges.  The other three vertices are
isolated in the nonloop support and each has at least two loop choices, giving
at least $2^3$ covers.

Suppose the matching number is two, and fix disjoint supports $E_1,E_2$ with
multiplicities $m_1,m_2$.  Let $u,v$ be the unmatched vertices.  Choosing one
actual point on each $E_i$ and one incident point at each of $u,v$ gives
\[
  m_1m_2d(u)d(v)
\]
distinct four-point covers.  They are minimal because no three-point
transversal exists.  This number is at least six unless
$m_1=m_2=1$ and $d(u)=d(v)=2$.

In that exceptional case, a second actual two-edge matching generates at
least four analogous covers.  The two families overlap in at most two.  If
the matchings share no actual point, a common cover must be their four-point
union.  If they share one actual point, the two remaining supports cannot be
disjoint.  When they meet in one vertex, the fourth point must cover a fixed
unmatched vertex of degree two; when they are parallel, no minimal common
cover exists.  Their union therefore contains at least six covers.

If the fixed matching is unique, every incidence at $u$ or $v$ is a loop.
Let $e_i$ be the unique actual point on $E_i$.  Write $E_1=\{a,b\}$ and
$E_2=\{c,d\}$, and choose incident points
$r_a\neq e_1$ and $r_b\neq e_1$ at $a$ and $b$.  Uniqueness forces
$r_a\neq r_b$ and keeps both supports inside the four matched vertices, so
$\{r_a,r_b,e_2\}$ covers those vertices.  An inclusion-minimal subcover
$K$ avoids $e_1$.  Combining $K$ with either loop at $u$ and either loop at
$v$ gives four minimal covers avoiding $e_1$, disjoint from the four covers
containing $e_1$.  Thus $b\geq6$.

The bounds in (2) now exhaust all possibilities:
\[
\begin{array}{c|c}
\text{case}&ab\\
\hline
b=2&\geq32\\
b=3&\geq36\\
b=4&\geq32\\
b\geq5,\ a=4&\geq64\\
b\geq5,\ a=5&\geq40\\
b\geq5,\ a=6&\geq36\\
b\geq5,\ a\geq7&\geq35
\end{array}
\]
which proves the proposition.
\end{proof}

\section{Layering predecessor certificate}
\label{app:layering}

In the following list, $B\leftarrow A$ means that $A$ belongs to the preceding
layer and $A\subsetneq B$.  A suffix $H$ means union with the full homogeneous
block.

\begin{description}[leftmargin=1.8em,style=nextline]
\item[$\mathcal A_1\leftarrow\mathcal A_0$]
$1\leftarrow\varnothing$;
$2\leftarrow\varnothing$;
$3\leftarrow\varnothing$;
$6\leftarrow\varnothing$;
$7\leftarrow\varnothing$;
$458H\leftarrow\varnothing$.

\item[$\mathcal A_2\leftarrow\mathcal A_1$]
$24\leftarrow2$; $15\leftarrow1$; $35\leftarrow3$;
$26\leftarrow2$; $47\leftarrow7$; $67\leftarrow6$;
$18\leftarrow1$; $38\leftarrow3$; $1346H\leftarrow1$;
$1237H\leftarrow1$; $4568H\leftarrow6$; $2578H\leftarrow2$.

\item[$\mathcal A_3\leftarrow\mathcal A_2$]
$124\leftarrow24$; $234\leftarrow24$; $145\leftarrow15$;
$345\leftarrow35$; $156\leftarrow15$; $356\leftarrow35$;
$267\leftarrow26$; $178\leftarrow18$; $378\leftarrow38$;
$12357H\leftarrow15$; $13467H\leftarrow47$;
$12358H\leftarrow15$; $12368H\leftarrow26$;
$13468H\leftarrow18$; $24568H\leftarrow24$;
$24578H\leftarrow24$; $45678H\leftarrow47$.

\item[$\mathcal A_4\leftarrow\mathcal A_3$]
$2345\leftarrow234$; $1247\leftarrow124$;
$3568\leftarrow356$; $1678\leftarrow178$;
$123567H\leftarrow156$; $134567H\leftarrow145$;
$123468H\leftarrow124$; $124568H\leftarrow124$;
$123578H\leftarrow178$; $134578H\leftarrow145$;
$234678H\leftarrow234$; $245678H\leftarrow267$.

\item[$\mathcal A_5\leftarrow\mathcal A_4$]
$23457\leftarrow2345$; $1234568H\leftarrow2345$;
$1234678H\leftarrow1247$; $1235678H\leftarrow3568$;
$1245678H\leftarrow1247$; $1345678H\leftarrow3568$.

\item[$\mathcal A_6\leftarrow\mathcal A_5$]
$12345678H\leftarrow23457$.
\end{description}

{\small
\bibliographystyle{cas-model2-names}
\bibliography{references}
}

\end{document}
